\documentclass[11pt]{article}
\usepackage[margin=1in]{geometry}
\usepackage{amsmath,amssymb,amsthm}
\usepackage[hidelinks]{hyperref}
\newtheorem{theorem}{Theorem}
\newtheorem{lemma}{Lemma}
\title{Exact zero-free gaps for every wide\\three-frequency progression of integer ratio}
\author{}
\date{28 November 2026}
\begin{document}
\maketitle
\vspace{-2em}
\begin{abstract}
For the real trigonometric spectrum \(\{\pm N,\pm(m+1)N,\pm(2m+1)N\}\),
we determine the supremal zero-free length for every integer \(m\ge2\).
Even \(m\) gives \(1/(2N)\). For every odd \(m\ge3\), a single
analytic branch gives the exact gap and an optimal positive measure
on six points. An interval certificate covers the entire parameter
range; the result has no unspecified large-\(m\) threshold.
\end{abstract}

\section{Statement}
Let \(M_m(N)\) be the supremal length of a zero-free interval
of a real trigonometric polynomial on \(\mathbb R/\mathbb Z\)
supported on \(\{\pm N,\pm(m+1)N,\pm(2m+1)N\}\).
Kozma and Oravecz asked for sharp lengths for wide arithmetic
progressions of frequencies \cite{KO}. ProveIt's spectral filter
chapter relates this question to the shortest support of a positive
annihilating measure \cite{ProveIt}. Here we settle the subfamily
with three frequencies and step \(mN\) for integer \(m\ge2\).
The result does not use RH or GRH.

For \(0\le\epsilon\le 1/3\) define
\[
 h_\epsilon(s)=
 \begin{cases}\sin(\epsilon s)/\epsilon,&\epsilon\ne0,\\s,&\epsilon=0,\end{cases}
 \quad C_\epsilon(s)=\cos((1+\epsilon)s),\qquad
 S_\epsilon(s)=\sin((2+\epsilon)s).
\]
Set \(L_\epsilon=\pi/(2\epsilon)\) for \(\epsilon>0\) and \(L_0=+\infty\).
The functions \(A,B,u,v,x\) used below are the unique choices with
\(x<0<u<v<L_\epsilon\) for which
\[
 F_\epsilon=h_\epsilon+A C_\epsilon+B S_\epsilon,\qquad
 F_\epsilon(u)=F_\epsilon(v)=F_\epsilon'(u)=F_\epsilon'(v)=F_\epsilon(x)=0
\]
and the positivity and dual sign conditions of Lemma~\ref{lem:cert} hold.
Primes on \(F\) mean differentiation in \(s\).

\begin{theorem}\label{thm:main}
For every positive integer \(N\) and integer \(m\ge2\),
\[
 NM_m(N)=
 \begin{cases}
  1/2,&m\ \text{even},\\[2pt]
  1/2-x(1/m)/(\pi m),&m\ \text{odd}.
 \end{cases}
\]
For odd \(m\) the supremum is not attained, even on an open interval.
On the fixed critical closed interval there is, up to positive scale,
a unique nonzero nonnegative polynomial, and a unique annihilating
probability measure, with exactly six atoms. Moreover, along odd \(m\),
\[
\begin{gathered}
 NM_m(N)=\frac12+\frac{c_1}{m}+\frac{c_2}{m^2}+O(m^{-3}),\\
 c_1=-x(0)/\pi \simeq 0.4313527,\qquad
 c_2=-x'(0)/\pi \simeq -0.2501628.
\end{gathered}
\]
\end{theorem}

\section{Certificate on the full continuum}
\begin{lemma}\label{lem:cert}
For every \(0\le\epsilon\le 1/3\), there is a solution of the five
equations above with \(-2<x<0<u<v<4\) such that:
\begin{enumerate}
\item \(F_\epsilon\ge0\) on \([x,L_\epsilon]\), with precisely the
zeros \(x,u,v\), and \(F_\epsilon'(x), F_\epsilon''(u),
F_\epsilon''(v)>0\).
\item \(C_\epsilon(x),C_\epsilon(u)>0>C_\epsilon(v)\).
Writing \(\rho(s)=S_\epsilon(s)/C_\epsilon(s)\) at these three points,
one has \(\rho(x)<\rho(v)<\rho(u)\).
\item The determinant with rows
\[
 \bigl(\cos(\epsilon s),\ \sin((1+\epsilon)s),\
                       \cos((2+\epsilon)s)\bigr),\qquad s=x,u,v,
\]
is positive.
\end{enumerate}
\end{lemma}

\begin{proof}
This is a finite interval certificate, replayed by
\texttt{node all\_odd.js}. The exact operations and their interpretation
are as follows. The bundled 588 rational cells can be replayed without calling any
floating-point proposals; the optional generation command rediscovers them.
The cells also give a reproducible definition of \(x(1/m)\)
without tracking a numerical optimizer.

\paragraph{Five simultaneous equations.}
Put \(z=(A,B,u,v,x)\) and
\(g=(F(u),F(v),F'(u),F'(v),F(x))\). Subdivide the closed parameter range
by bisection. For a parameter cell \(E\), a floating-point routine proposes
a center \(z_0\) and preconditioner \(R\), both converted to rational
grid points. It makes no decisions about validity. The checker
encloses \(R g(E,z_0)\) and \(I-R D_zg(E,Z)\) on a common-radius box
\(Z=z_0+[-r,r]^5\). For each row it verifies
\[
 q_i:=\sum_j\sup_Z |(I-RD_zg)_{ij}|<1,\qquad
 \sup_E |(Rg(E,z_0))_i|+rq_i<r.
\]
For every parameter in the cell the map \(z\mapsto z-Rg(z)\)
therefore sends \(Z\) into its interior and is a contraction in
the maximum norm. The bound on \(I-RD_zg\) implies \(R\) and
\(D_zg\) are nonsingular, so the fixed point is a zero of \(g\).
The arithmetic subdivides any cell whose verification fails;
all accepted cells together cover \([0,1/3]\).

\paragraph{Positivity on the whole physical interval.}
The checker verifies \(-2<x<0<u<v<4\) throughout each root box.
Extend each interval enclosure of \(x,u,v\) by \(1/9\)
(rounded down to the fixed grid); the enlarged intervals are
disjoint. It verifies \(F'>0\) on the first and \(F''>0\) on the
other two. These checks give positivity near the simple endpoint and
away from the centers of the two double zeros. Between these
neighborhoods and up to the end of the compact part, interval
bisection verifies \(F>0\). In a cell \(E\) with lower endpoint
\(\epsilon_-\), this compact check ends at
\(\min(7,11/(7\epsilon_-))\), rounded up unless the cap is 7, with the
zero denominator interpreted as \(+\infty\). Since \(\pi<22/7\), that
covers \([x,\min(7,L_\epsilon)]\). For \(7\le s\le L_\epsilon\),
\[
 h_\epsilon(s)\ge 2s/\pi\ge14/\pi>49/11>|A|+|B|;
\]
the last strict bound is checked on every box. The remaining tests are
the displayed cosine signs, slope inequalities and odd determinant.

\paragraph{Exact trigonometric enclosures.}
The arithmetic represents intervals on a \(10^{-23}\) grid and rounds
outward using integer floor and ceiling. Sine and cosine are evaluated
at the interval midpoint divided by eight, using their Taylor
polynomials through degrees 43 and 42, then three double-angle steps.
The bound \(3^{44}/44!<10^{-23}\) controls both remainders when the
reduced argument has absolute value at most three. The polynomial for
\(\sin y/y=\int_0^1\cos(ty)\,dt\) is bounded the same way. Input
interval radii are added by the Lipschitz bounds \(1,1,1/2\),
respectively. The replay checks all Taylor domains. Only rational
encapsulations of the numerical proposals enter the proof.
\end{proof}

\section{From certificates to sharp lengths and uniqueness}
At a fixed parameter, the signs and ordered slopes give positive
weights \(\lambda_x,\lambda_u,\lambda_v\) that cancel both \(C\) and
\(S\). Explicitly, select \(0<q<1\) with
\(\rho(v)=q\rho(x)+(1-q)\rho(u)\) and take
\[
 \lambda_x=\frac{q}{C(x)},\qquad
 \lambda_u=\frac{1-q}{C(u)},\qquad
 \lambda_v=-\frac1{C(v)}.
\]
Since \(F\) vanishes at all three points, these weights also cancel
\(h\). This proves uniqueness of the certified solution for a given
parameter. Indeed, if two endpoints satisfied \(x_1<x_2<0\), the first
polynomial would be nonnegative on the second contact set and positive
at \(x_2\), contrary to the second weights. For equal endpoints,
the same argument forces the second contact set to be the first.
The two distinct slopes at \(u,v\) determine \(A,B\) from the
two equations \(F(u)=F(v)=0\).

The fixed point lies in the interior of its box, and its Jacobian is
nonsingular. The analytic implicit function theorem thus gives an
analytic branch on each cell and through its endpoints. At shared
endpoints these branches coincide by the uniqueness just shown;
local uniqueness from the implicit function theorem glues them.
In particular there is an analytic germ through \(\epsilon=0\).

Dilation reduces the interval problem to \(N=1\). Write the angular
coordinate as \(\theta=2\pi t\). For odd \(m\), set
\(s=m(\pi/2-\theta)\) and \(\epsilon=1/m\). The three cosine modes
become \(h_\epsilon/m\), \((-1)^{(m+1)/2} C_\epsilon\), and
\(-S_\epsilon\). The critical half-angle is
\(\theta_*=\pi/2-x/m\). Place the three weights in symmetric pairs at
angles \(\pm(\pi/2-s/m)\), \(s=x,u,v\), and normalize to total mass one.
The cosine modes cancel as above; the sine modes cancel by reflection.
Any polynomial with no zero on an open interval of angular length
\(2\theta_*\) can be centered there by translation and made positive
on the interior by a sign change. Its endpoint values are nonnegative,
so integrating over the six atoms, four interior, is strictly
positive. This contradicts annihilation; larger intervals also fail.

For \(x<a<0\), the function \(F_\epsilon+\delta h_\epsilon\) is
strictly positive on \([a,L_\epsilon]\) for small positive \(\delta\):
on \([a,0]\) the unperturbed function has a positive minimum, and on
\((0,L_\epsilon]\) the perturbation is positive. It is a combination
of the allowed modes. Its zero-free intervals tend to the critical
length. Division by \(2\pi\) and dilation by \(N\) give the odd case.
For even \(m\), all three frequencies after dilation are odd. A shift
by \(1/2\) changes the sign of every polynomial, so any open interval
of length greater than \(1/2\) contains a zero. The single cosine
has a zero-free open interval of length \(1/2\).

For odd \(m\), any nonnegative polynomial on the critical interval
vanishes at all six points by the positive weights. The even part
is proportional to the certified function, since the two
\(C,S\) evaluations at \(u,v\) are independent. The odd part is zero
by Lemma~\ref{lem:cert}(3), which is precisely the evaluation
determinant of the sine modes up to fixed column signs.
The scalar is nonnegative. Conversely, an annihilating probability
measure is supported on the zeros of the nonnegative certified
polynomial, so is supported on the same six points. The sine
determinant forces equal masses in each pair. The cosine equations
have rank two in the three pair totals by the two slope evaluations.
Normalization fixes the unique pair totals, all positive.

\section{Expansion and reproducibility}
The germ gives \(x(1/m)=x(0)+x'(0)/m+O(m^{-2})\).
To calculate \(x'(0)\), set \(H(s)=-A(0)s\sin s+B(0)s\cos(2s)\),
the partial derivative of \(F_\epsilon\) holding the coefficients
fixed at \(\epsilon=0\). The equations at the two stationary
contacts give
\[
 A'(0)\cos s+B'(0)\sin(2s)=-H(s),\qquad s=u(0),v(0).
\]
At the simple endpoint the derivative is then
\[
 x'(0)=-\frac{A'(0)\cos x+B'(0)\sin(2x)+H(x)}{F_0'(x)},
 \qquad x=x(0).
\]
The separate \texttt{limit.js} encloses the base solution and gives
\(-1.355137<x(0)<-1.355132\) and
\(0.7859003465<x'(0)<0.7859187030\). It brackets the stationary
contacts using \(T(s,B)=-s/\cos s-2B\sin s\), and solves
\(T(u(B),B)=T(v(B),B)\) by strict monotonicity in \(B\).
The displayed decimal asymptotic coefficients are approximations;
the exact constants are the germ expressions in Theorem~\ref{thm:main}.
The full-parameter certificate needs no outputs from this auxiliary
limiting check. For example, numerical values of \(NM_m(N)\) for
\(m=3,5,7,9\) are \(0.618557753484\), \(0.576776857359\),
\(0.556692411116\), and \(0.544918838243\). The \(m=3\) specialization
agrees with the separate Chebyshev-polynomial certificate.

\begin{thebibliography}{9}
\bibitem{KO} G. Kozma and F. Oravecz,
\emph{On the gaps between zeros of trigonometric polynomials},
Real Anal. Exchange 28 (2002/03), no.~2, 447--454.
\href{https://arxiv.org/abs/math/0601021}{arXiv:math/0601021},
Question~1.
\bibitem{ProveIt}
ProveIt,
\href{https://github.com/VladimirReshetnikov/ProveIt/blob/3c4d27a69c98d09319a51e41a19fca05bb80a3e2/Analysis/GeneralizedRiemannHypothesis/Article/sections/fil.tex}{GRH article, spectral-filter chapter},
revision \texttt{3c4d27a69c98}.
\end{thebibliography}
\end{document}
